% Activate the following line by filling in the right side. If for example the name of the root file is Main.tex, write
% "...root = Main.tex" if the chapter file is in the same directory, and "...root = ../Main.tex" if the chapter is in a subdirectory.
 
%!TEX root =  ../main.tex

\chapter[GoogleAIMode]{Ideas Stolen From Google AI}

\section{Introduction}



\subsection{The Metric and Spacetime Interval}

To track a continuous navigation path across your 3D spatial field, your engine must continuously evaluate the localized spacetime interval \(ds^{2}\) using the coordinate position components \(dx^{\mu }\) and \(dx^{\nu }\), weighted by your localized metric tensor \(g_{\mu \nu }\)

\[
ds^2 = \sum_{\mu=0}^{3} \sum_{\nu=0}^{3} g_{\mu\nu} dx^\mu dx^\nu
\]


\subsection{The Christoffel Symbols (The Acceleration Field)}

Your background C++ CSym engine calculates the local gravitational or inertial acceleration using the Christoffel symbols of the second kind \(\Gamma _{\mu \nu }^{\lambda }\). This evaluates the spatial derivatives of your metric tensor components (\(g_{\nu \sigma }\) and \(g_{\mu \sigma }\)):

\[
\Gamma^\lambda_{\mu\nu} = \frac{1}{2} g^{\lambda\sigma} \left( \frac{\partial g_{\nu\sigma}}{\partial x^\mu} + \frac{\partial g_{\mu\sigma}}{\partial x^\nu} - \frac{\partial g_{\mu\nu}}{\partial x^\sigma} \right)
\]


\subsection{The Geodesic Differential Equation (The Real-Time Path)}

To render a continuous line in your visionOS volume, your background actor takes the Christoffel outputs and solves this second-order differential equation with respect to your proper affine parameter \(\tau \). This calculates the exact curving flight path across space-time:

\[
\frac{d^2 x^\lambda}{d\tau^2} + \sum_{\mu=0}^{3} \sum_{\nu=0}^{3} \Gamma^\lambda_{\mu\nu} \frac{dx^\mu}{d\tau} \frac{dx^\nu}{d\tau} = 0
\]


\subsection{Topological Cellular Sewing (The G-Map Boundary Constraint)}

When a geodesic trajectory hits the edge of a cell inside your Swift G-Map, your system sews the boundary face using your \(\alpha _{i}\) involutions. To guarantee the macro-topology remains a seamless, watertight manifold when executing cell contraction or removal, the structural lookahead must satisfy this algebraic condition:


\[
\alpha_i(\alpha_j(d)) = \alpha_j(\alpha_i(d)) \quad \text{for all } |i - j| \ge 2
\]



\chapter{From Sailing Tacks to Spatial Cosmological Foliations}
\label{ch:sailing_to_cosmology}

\section{The Architectural Extension}
The foundational mechanism of our geometry workbench relies on a core abstraction: \textit{unfolding space by reflections} rather than turning the path \cite{damiand2014}. In our initial two-dimensional formulation for a sailing tacking optimizer, the configuration space was partitioned into continuous level sets defined by a smooth scalar field $s(x,y) = c$. 

To lift this system into a real-time, interactive three-dimensional environment for spatial computing platforms (such as Apple visionOS), we must generalize the discrete alternating tack states into a dynamic, multi-dimensional metric field. The topology is managed natively via a specialized Generalized Map (\GMap) allocation structure, while local geometric constraints are evaluated through an unknotted, watertight manifold kernel.

\section{The Mathematical Formulation}

To achieve fluid, real-time spatial path rendering at high update loops ($90\text{Hz}$--$120\text{Hz}$), the physical computation layer is split from the visual instrumentation bus. The underlying continuous physics are governed by three primary geometric equations.

\subsection{The Local Spacetime Interval}
The continuous navigation path across an arbitrary topological coordinate field is tracked via the localized interval $ds^2$. This metric is a function of coordinate displacement vectors $dx^\mu$ and $dx^\nu$, scaled by the localized metric tensor $g_{\mu\nu}$:

\[
ds^2 = \sum_{\mu=0}^{3} \sum_{\nu=0}^{3} g_{\mu\nu} dx^\mu dx^\nu
\]

\subsection{The Acceleration Connection Field}
The symbolic parsing engine evaluates localized geometric acceleration tensors using Christoffel symbols of the second kind, denoted as $\Gamma^\lambda_{\mu\nu}$. This calculation handles the partial spatial derivatives of the underlying metric components:

\[
\Gamma^\lambda_{\mu\nu} = \frac{1}{2} g^{\lambda\sigma} \left( \frac{\partial g_{\nu\sigma}}{\partial x^\mu} + \frac{\partial g_{\mu\sigma}}{\partial x^\nu} - \frac{\partial g_{\mu\nu}}{\partial x^\sigma} \right)
\]

\subsection{The Real-Time Geodesic Path}
The trajectory of an actively simulated particle (or navigation vector) is calculated by solving the second-order differential equation with respect to the proper affine parameter $\tau$. This guides the real-time path generation thread:

\[
\frac{d^2 x^\lambda}{d\tau^2} + \sum_{\mu=0}^{3} \sum_{\nu=0}^{3} \Gamma^\lambda_{\mu\nu} \frac{dx^\mu}{d\tau} \frac{dx^\nu}{d\tau} = 0
\]

\section{Topological Cell Constraints/Topological Cellular Sewing (The G-Map Boundary Constraint)}
When a continuous geodesic trajectory reaches a boundary interface within the underlying mesh, the structural layout is updated via localized $\alpha_i$ involutions. To guarantee a valid, watertight boundary manifold during dynamic cell contraction or removal operations, the combinatorial map lookahead must strictly satisfy the following execution identity:

\[
\alpha_i(\alpha_j(d)) = \alpha_j(\alpha_i(d)) \quad \text{for all } |i - j| \geq 2
\]



\chapter{From Sailing Tacks to Spatial Cosmological Foliations}
\label{ch:sailing_to_cosmology}

\section{The Architectural Extension}
The foundational mechanism of our geometry workbench relies on a core abstraction: \textit{unfolding space by reflections} rather than turning the path \cite{damiand2014}. In our initial two-dimensional formulation for a sailing tacking optimizer, the configuration space was partitioned into continuous level sets defined by a smooth scalar field $s(x,y) = c$. 

To lift this system into a real-time, interactive three-dimensional environment for spatial computing platforms (such as Apple visionOS), we must generalize the discrete alternating tack states into a dynamic, multi-dimensional metric field. The topology is managed natively via a specialized Generalized Map (\GMap) allocation structure, while local geometric constraints are evaluated through an unknotted, watertight manifold kernel.

\section{The Mathematical Formulation}

To achieve fluid, real-time spatial path rendering at high update loops ($90\text{Hz}$--$120\text{Hz}$), the physical computation layer is split from the visual instrumentation bus. The underlying continuous physics are governed by three primary geometric equations.

\subsection{The Local Spacetime Interval}
The continuous navigation path across an arbitrary topological coordinate field is tracked via the localized interval $ds^2$. This metric is a function of coordinate displacement vectors $dx^\mu$ and $dx^\nu$, scaled by the localized metric tensor $g_{\mu\nu}$:

\[
ds^2 = \sum_{\mu=0}^{3} \sum_{\nu=0}^{3} g_{\mu\nu} dx^\mu dx^\nu
\]

\subsection{The Acceleration Connection Field}
The symbolic parsing engine evaluates localized geometric acceleration tensors using Christoffel symbols of the second kind, denoted as $\Gamma^\lambda_{\mu\nu}$. This calculation handles the partial spatial derivatives of the underlying metric components:

\[
\Gamma^\lambda_{\mu\nu} = \frac{1}{2} g^{\lambda\sigma} \left( \frac{\partial g_{\nu\sigma}}{\partial x^\mu} + \frac{\partial g_{\mu\sigma}}{\partial x^\nu} - \frac{\partial g_{\mu\nu}}{\partial x^\sigma} \right)
\]

\subsection{The Real-Time Geodesic Path}
The trajectory of an actively simulated particle (or navigation vector) is calculated by solving the second-order differential equation with respect to the proper affine parameter $\tau$. This guides the real-time path generation thread:

\[
\frac{d^2 x^\lambda}{d\tau^2} + \sum_{\mu=0}^{3} \sum_{\nu=0}^{3} \Gamma^\lambda_{\mu\nu} \frac{dx^\mu}{d\tau} \frac{dx^\nu}{d\tau} = 0
\]

\section{Topological Cell Constraints}
When a continuous geodesic trajectory reaches a boundary interface within the underlying mesh, the structural layout is updated via localized $\alpha_i$ involutions. To guarantee a valid, watertight boundary manifold during dynamic cell contraction or removal operations, the combinatorial map lookahead must strictly satisfy the following execution identity:

\[
\alpha_i(\alpha_j(d)) = \alpha_j(\alpha_i(d)) \quad \text{for all } |i - j| \geq 2
\]



